% GPT-6 Astra Pro; effort not exposed; /root/replacement; 22 November 2026.
\section[HD-F1-R26: the exact rank-two character gate]{HD-F1-R26: the remaining adjacent endpoints share one of two explicit rank-two character tensors}
\label{sec:followup-rm-rank-two-character-gate}

Retain the genuine real-quadratic fourfold, compatible
polarization and balanced vector
$v_q=v_1v_2$, $v_i=1-qA_i$, from R07--R19.
We first classify every rational Hodge vector
$z\in\bigoplus_pH^{p,p}(X)\cap H^{\rm even}(X,\mathbf Q)$
for which
\[
                  r(z):=\operatorname{rank}
                     \bigl(K\xrightarrow{k\mapsto D_kz}HH_{-2}(X)\bigr)
                                      \le3,\qquad K=\ker c^2_{v_q}.
 \tag{F1R26.1}
\]
Its rank is zero or two. At rank zero it is a rational
multiple of $v_q$. At rank two it has exactly one of
the two forms (F1R26.15)--(F1R26.16) below, with
all tensor, field-involution and trace maps specified.

Applied to every remaining actual adjacent object
of R21 at endpoint ranks $(3,2)$, this proves that
both endpoint whole characters have rank exactly two
in (F1R26.1). Their nonbalanced tensors are identical,
and their balanced coefficients differ by one.
This constrains all Chern-character components
together, while retaining the actual surface
defects and coupling of R22--R25.

\paragraph{1. Reality removes the off-middle block degrees.}
Use the actual cyclic exterior isomorphism of R19,
\[
 \Phi=\Phi_1\widehat\otimes\Phi_2:
       \bigwedge^4(U_1\oplus U_2)\xrightarrow{\sim}HH_0(X),
 \qquad
 \Phi_i(a)=D_a(x_{i1}x_{i2}).
 \tag{F1R26.2}
\]
It carries multiplication by $K_i$ to the
characteristic action. Write
$w=\Phi^{-1}(z)=\sum_{j=0}^4w_j$, where
$w_j\in\bigwedge^jU_1\otimes\bigwedge^{4-j}U_2$.
The argument in R19 gives rank at least five
from a nonzero $w_0$ or $w_4$: restrict respectively
to $K_1$ or $K_2$ and project to block degrees
$(2,4)$ or $(4,2)$. Consequently both vanish
under (F1R26.1).

For completeness, the precise next rank step is
stronger than the one used in R19.
For a nonzero $a\in U_i$, the kernel of
$k\mapsto ka$ on $\bigwedge^2U_i$ is
$a\wedge U_i$, of dimension three.
The nondegenerate alternating form defining $K_i$
is nonzero on that subspace. Thus its intersection
with $K_i$ has dimension two and
\[
                    \operatorname{rank}(K_i\to\bigwedge^3U_i,
                                              \ k\mapsto ka)=3.
 \tag{F1R26.3}
\]
For independent $a,b\in U_i$ the two kernels on
the full exterior square intersect in
$\mathbf C(a\wedge b)$. Hence a tensor
$w_1\in U_1\otimes\bigwedge^3U_2$ of tensor rank
at least two makes its $K_1$ action have rank at
least four. The same assertion holds for $w_3$
and $K_2$.

Reality of $z$ is used here through the actual
complex-conjugation map on forms. On the $i$th
block, exterior degree $j$ in (F1R26.2) has
Hochschild homological degree $2-j$. Complex
conjugation reverses that degree. It therefore
exchanges the nonvanishing of $w_1$ and $w_3$,
without exchanging the RM eigenspaces.
Under (F1R26.1), any nonzero pair must consequently
have the form
\[
                   w_1=a\otimes\beta_2,\qquad
                   w_3=\beta_1\otimes b
 \tag{F1R26.4}
\]
with all four factors nonzero. Their two projected
action images in
$\bigwedge^3U_1\otimes\bigwedge^3U_2$ have dimension
three and are respectively
\[
              (K_1a)\otimes\mathbf C\beta_2,\qquad
                      \mathbf C\beta_1\otimes(K_2b).
 \tag{F1R26.5}
\]
Their intersection has dimension at most one:
a tensor in it must lie in
$\mathbf C(\beta_1\otimes\beta_2)$.
Their sum therefore has dimension at least five,
contradicting (F1R26.1). We have proved
\[
                        w=w_2\in
                             \bigwedge^2U_1\otimes\bigwedge^2U_2.
 \tag{F1R26.6}
\]
This step used complex conjugation only in its
displayed Hodge-degree action.

\paragraph{2. The actual RM projectors and the field-defined pairings.}
Let $F=\mathbf Q(\sqrt{\Delta_0})$ and choose the
actual endomorphism
$j=m(1+\sqrt{\Delta_0})\in\operatorname{End}(X)$
as in R19. Put
$\alpha_1=m(1+\sqrt{\Delta_0})$ and
$\alpha_2=m(1-\sqrt{\Delta_0})$.
On $H^1(X,\mathbf Q)\otimes_{\mathbf Q}F$ the
explicit polynomial projectors are
\[
 e_1=\frac{j^*-\alpha_2}{\alpha_1-\alpha_2},
       \qquad
 e_2=\frac{j^*-\alpha_1}{\alpha_2-\alpha_1},
       \qquad H_i=e_i(H^1(X,\mathbf Q)\otimes F).
 \tag{F1R26.7}
\]
They satisfy $e_1+e_2=1$, $e_i^2=e_i$ and
$e_1e_2=0$; each $H_i$ has dimension four over $F$.
Exterior product gives the actual isomorphism
$\bigwedge^*H_1\widehat\otimes_F\bigwedge^*H_2
\to H^*(X,\mathbf Q)\otimes F$.
Set
\[
 S_i=\bigwedge\nolimits^{\rm even}H_i,\qquad
 B_i(s,t)=\frac{[s^\vee t]_4}{A_i},\qquad
 s^\vee=\sum_{p=0}^2(-1)^p s_{2p}.
 \tag{F1R26.8}
\]
The quotient in this formula takes the coefficient
in the one-dimensional space $\bigwedge^4H_i=F A_i$.
It defines a nondegenerate symmetric $F$-bilinear
pairing on the eight-dimensional $S_i$.
The compatible polarization makes
$\theta_i=(\bigwedge^2e_i)\theta$ and
$A_i=\theta_i^2/2$ the stated classes over $F$.
Their conjugate sums are rational. In particular
the rationality of $v_q$ makes $q$ rational:
its degree-four term is $-q(A_1+A_2)$, and
$A_1+A_2$ is a nonzero rational vector.

Write
\[
           V_i=v_i^{\perp B_i}\subset S_i,\qquad
           B_i(v_i,v_i)=-2q,\qquad
           S_i=Fv_i\oplus V_i.
 \tag{F1R26.9}
\]
After the chosen embedding into $\mathbf C$, put
\[
 S_i^0=H^{0,0}_i\oplus H^{1,1}_i\oplus H^{2,2}_i.
 \tag{F1R26.10}
\]
It has dimension six. The map
$\Phi_i:\bigwedge^2U_i\to S_i^0$ is the actual
isomorphism from (F1R26.2).
With the basis order $(t_1,t_2,w_1,w_2)$ and
$A_i=x_1x_2y_1y_2$, direct exterior multiplication
gives, for $a,b\in\bigwedge^2U_i$,
\[
          B_i(\Phi_i(a),\Phi_i(b))
                   =-[ab]_{t_1t_2w_1w_2}.
 \tag{F1R26.11}
\]
For example $\Phi_i(t_1t_2)=-1$,
$\Phi_i(w_1w_2)=A_i$,
$\Phi_i(t_1w_1)=x_2y_1$ and
$\Phi_i(t_2w_2)=-x_1y_2$; all other complementary
pairs follow with their exterior signs.
Also $\Phi_i(\xi_i)=v_i$, for
$\xi_i=-t_1t_2-qw_1w_2$.
Thus $\Phi_i(K_i)=V_i\cap S_i^0$.
The pairing orthogonal complement of $S_i^0$
in $S_i\otimes_F\mathbf C$ is
$H^{2,0}_i\oplus H^{0,2}_i$, which is itself
orthogonal to $v_i$. Consequently
\[
 V_i\otimes_F\mathbf C
      =(V_i\cap S_i^0)\oplus H^{2,0}_i\oplus H^{0,2}_i.
 \tag{F1R26.12}
\]

By (F1R26.6), the rational vector $z$ belongs,
after scalar extension, to $S_1^0\otimes S_2^0$.
Under exterior product it is also an $F$-tensor
in $S_1\otimes_FS_2$.
Contract it with $B_1$ or $B_2$ on $V_1$ or
$V_2$, obtaining actual $F$-linear maps
$C_1(z):V_1\to S_2$ and $C_2(z):V_2\to S_1$.
The two extra directions in (F1R26.12) contract
$z$ to zero. Equation (F1R26.11) shows that
the restricted contraction is the $K_i$ action,
multiplied by the fixed nonzero vector
$-\Phi_i(t_1t_2w_1w_2)$. In particular their ranks
are equal. The two $K_i$ images have distinct
block degrees $(4,2)$ and $(2,4)$, so
\[
                          r(z)=\operatorname{rank}C_1(z)
                                      +\operatorname{rank}C_2(z).
 \tag{F1R26.13}
\]

The required field comparison now takes place
on rational cohomology with scalar extension.
Let a bar denote the semilinear map
$1\otimes\tau$, where $\tau$ is the nontrivial
automorphism of $F/\mathbf Q$.
It sends $e_1,H_1,A_1,v_1,B_1$ to their
index-two counterparts, and conversely, and
fixes $z$. The resulting maps intertwine
$C_1(z)$ and $C_2(z)$ semilinearly.
Their ranks over $F$, and hence over $\mathbf C$,
are equal. This is an explicit scalar-extension
map; it makes no assertion that $\tau$ preserves
the Hodge filtration. The ranks were compared
to the actual Hodge action beforehand by
(F1R26.11)--(F1R26.13).
It follows that $r(z)$ is even under (F1R26.1).
Thus it is zero or two.

\paragraph{3. The complete two-case tensor classification.}
The $F$-orthogonal splittings (F1R26.9) write
$z$ uniquely as
\[
 z=a\,v_1v_2+b_1v_2+v_1b_2+M,\qquad
          b_i\in V_i,\quad M\in V_1\otimes_FV_2.
 \tag{F1R26.14}
\]
The displayed tensors mean their exterior-product
images in actual cohomology. Their uniqueness
and the field comparison imply
$a\in\mathbf Q$, $b_2=\overline{b_1}$, and
$\overline M=M$ with the tensor factors
interchanged. After embedding in $\mathbf C$,
each $b_i$ lies in $S_i^0$ and
$M$ lies in $(V_1\cap S_1^0)\otimes(V_2\cap S_2^0)$.

If $r(z)=0$, both contractions in (F1R26.13)
vanish. Nondegeneracy on $V_i$ then gives
$b_1=b_2=M=0$, so $z=a v_q$.
Suppose now $r(z)=2$. Each contraction has rank one.
When $M=0$ this forces $b_1\ne0$ and gives the
first exact form
\[
             z=a\,v_q+b\,v_2+v_1\overline b,\qquad
             a\in\mathbf Q,\quad
             0\ne b\in V_1\cap S_1^0
                       \text{ defined over }F,\quad
                       \overline b\in S_2^0.
 \tag{F1R26.15}
\]

When $M\ne0$, its contraction rank is one, so
write $M=u_1\otimes u_2$ with nonzero $u_i\in V_i$
over $F$. Rank one of the full first contraction
forces $b_1=s u_1$ for some $s\in F$:
otherwise its two coefficient functionals into
$Fv_2\oplus Fu_2$ would be independent.
Likewise $b_2=t u_2$ for $t\in F$.
The actual tensor identity is
\[
 z=(a-st)v_1v_2+(u_1+t v_1)(u_2+s v_2).
\]
Moreover $b_1\otimes b_2=st\,M$, so uniqueness
of that scalar and the field involution give
$st\in\mathbf Q$. Put $a_0=a-st\in\mathbf Q$.
The rank-one $F$-tensor $z-a_0v_q$ is fixed by
the factor-exchanging involution. Choose one
nonzero factor $p\in S_1$ over $F$. Its second
factor must be $c\,\overline p$ for $c\in F^*$;
applying the involution again gives $c=\overline c$.
Thus the second exact form is
\[
             z=a_0v_q+c\,p\,\overline p,\qquad
             a_0\in\mathbf Q,\quad c\in\mathbf Q^*,\quad
             p\in S_1^0\text{ defined over }F,\quad
             \overline p\in S_2^0,\quad
             p\notin Fv_1.
 \tag{F1R26.16}
\]
The factors lie in their indicated Hodge-zero
spaces because the nonzero decomposable tensor
$z-a_0v_q$ lies in $S_1^0\otimes S_2^0$.
Conversely (F1R26.15) has one nonzero contraction
in each block, and so does (F1R26.16), since
the projections of $p,\overline p$ to $V_1,V_2$
are nonzero. Both therefore have rank exactly two.
The two cases are distinguished by whether
$M$ in the unique decomposition (F1R26.14)
vanishes. No factor in these cohomological
identities is declared to be the character
of a sheaf.

\paragraph{4. The strict whole-character gate for actual endpoints.}
Consider an actual remaining object of R21
at ranks $(r_A,r_T)=(3,2)$, with its determinant
twist already performed. The common-quotient
branch has $h=2$ by R22; the common-image
branch has $h=2$ or $3$ by R23. The actual
endpoint Yoneda memberships give
\[
 \operatorname{ev}_A^2(K)\subset
                   \delta\operatorname{Hom}_G(A,T),
       \qquad
 \operatorname{ev}_T^2(K)\subset
                   \operatorname{Hom}_G(A,T)[2]\delta.
 \tag{F1R26.17}
\]
Taking the ordinary full trace shows
$r(\operatorname{ch}(A)),r(\operatorname{ch}(T))\le h\le3$.
The actual $K_0$ equality
$\operatorname{ch}(A)-\operatorname{ch}(T)=v_q$
makes their restricted character actions equal.
Hence both have the same rank, zero or two.

Rank zero is excluded for these very endpoints.
It would give
$\operatorname{ch}(A)=3v_q$ and
$\operatorname{ch}(T)=2v_q$.
The proved saturated-filtration and actual
double-dual Hom estimate (F1R18.20)--(F1R18.25),
using the complete DEG\_L S-22 effective-cycle
argument in its all-one-line-hull case, would give
\[
 d:=\dim\operatorname{End}_G(A)+\dim\operatorname{End}_G(T)
                         \le5+2=7.
 \tag{F1R26.18}
\]
The unrestricted tangent/endpoint inequality of
R20 retains this same $h$ and gives
\[
 c+b+\frac{k}{2}\le d+1+h\le11,\qquad
 c=\chi_G(T,A),\quad b=\operatorname{ext}^1_G(A,T)\ge0.
 \tag{F1R26.19}
\]
For the two balanced endpoint characters the
actual polarized invariant Euler identity of
DEG\_L S-21, already used in R18, gives
$c=6k$, where $k=\chi_G(E,E)$ is a positive
even integer. The left side of (F1R26.19) is
therefore at least $13$, a contradiction.
Thus
\[
           r(\operatorname{ch}(A))
             =r(\operatorname{ch}(T))=2.
 \tag{F1R26.20}
\]

Apply the unique decomposition (F1R26.14) to
both endpoint characters. Their difference
$v_q$ changes only its first coefficient.
Consequently they have one of the two
complete patterns
\[
 \begin{aligned}
 \operatorname{ch}(T)&=a v_q+bv_2+v_1\overline b,&
 \operatorname{ch}(A)&=(a+1)v_q+bv_2+v_1\overline b,\\
 \text{or}\qquad
 \operatorname{ch}(T)&=a_0v_q+c p\overline p,&
 \operatorname{ch}(A)&=(a_0+1)v_q+c p\overline p,
 \end{aligned}
 \tag{F1R26.21}
\]
with the exact nonzero and rationality conditions
in (F1R26.15)--(F1R26.16), respectively.
The formulas are identities for the existing
whole characters, not proposed substitutes for
the actual sheaves or for $\delta$.

\paragraph{Framework, checks and remaining scope.}
The characteristic map, full trace and exact
endpoint Yoneda kernels are the same
Buchweitz--Flenner maps checked in R07 and R19.
The polynomial RM projectors (F1R26.7),
exterior product, normalized Mukai pairings
(F1R26.8), and scalar-extension involution
give the complete comparison used for the
new rank restriction. An independent exact
four-variable exterior calculation checks
all thirty-six basis pairings in (F1R26.11);
the proof above also gives their signs directly.
The earlier rational $B$-field is restored
by the actual invertible shear of R07,
$c^2_{e^{-B}z}=e^{-B}c^2_z\bigwedge^2s_B$,
so the same rank statement applies to its
specified unnormalized characteristic map.

This result eliminates every endpoint character
outside (F1R26.21), including every candidate
whose restricted character action has rank
three. It imposes the exact off-middle block
vanishing and common tensor on the surviving
surface-cancellation branches. It neither
constructs their sheaves nor removes their
surface Ext contributions. The actual finite
action, full character and connecting morphism
remain the ones in R21--R25.
